% ====================================================================
% DATEI: content/12_collapse_rebirth_de.tex
% Kollaps und Wiedergeburt von Kontinua
% Core 1.3 — kanonische Version
% ====================================================================

\section{Kollaps und Wiedergeburt von Kontinua}
\label{sec:collapse-rebirth}

Dieser Abschnitt entwickelt die strukturelle Theorie von Kollaps und
Wiedergeburt in der Ontology of Continua (OC). Er verfeinert die allgemeine
Todesbedingung \(\Omega(K)=\emptyset\) und \(k(K,t)\to 0\) aus
Abschnitt~\ref{sec:results} zu expliziten Kriterien, dynamischen Signaturen und
restriktierten Szenarien für die Entstehung neuer Kontinua nach einem Kollaps.

Es werden keine neuen Axiome eingeführt; die Diskussion fasst die
Folgerungen des allgemeinen Rahmens aus Abschnitt~\ref{sec:model}, der
strukturellen Ergebnisse aus Abschnitt~\ref{sec:results}, der erweiterten
Randbehandlung aus Abschnitt~\ref{sec:boundary-extended}, der
Schwellwertlandschaft aus Abschnitt~\ref{sec:thresholds-extended}, der
vollständigen Ebenenhierarchie aus Abschnitt~\ref{sec:klevels-full} und der
Operatorfamilie aus Abschnitt~\ref{sec:operators-full} zusammen.

\subsection{Formale Kriterien für Kollaps}

Sei
\[
  K(t) =
  \big(
    \Omega(t), A(t), P(t), J(t),
    \Theta(t), \partial\Omega(t), C(t), k(t)
  \big)
\]
ein Kontinuum, eingebettet in einen Raum \(M\). Der strukturelle Kollapsbegriff
muss vorübergehende Exkursionen zur Grenze von einem echten Ende des
Kontinuums unterscheiden.

\paragraph{Definition 12.1 (Kollaps).}
Ein Kontinuum \(K\) \emph{kollabiert} zum Zeitpunkt \(t^{\ast}\), wenn folgende
Bedingungen erfüllt sind:
\begin{enumerate}
  \item der zulässige Zustandsraum wird leer:
        \[
          \Omega\big(K(t^{\ast})\big) = \emptyset;
        \]
  \item die Continuumness fällt auf Null:
        \[
          \lim_{t \nearrow t^{\ast}} k\big(K,t\big) = 0;
        \]
  \item mindestens ein Existenz- oder Stabilitätsschwellwert ist für alle
        Kandidatenkonfigurationen verletzt:
        \[
          \forall s \in \Omega(M):
          \quad
          \exists \theta \in \Theta_{\mathrm{exist}} \cup \Theta_{\mathrm{stab}}
          \ \text{so dass}\ \theta(s) > 0.
        \]
\end{enumerate}

Diese Bedingungen entsprechen den äquivalenten Charakterisierungen des Todes aus
Satz~3 und Satz~9 in Abschnitt~\ref{sec:results}. Insbesondere fallen das
Verschwinden von \(\Omega\) und \(k\) mit dem Verschwinden des maximalen
strukturell stabilen Zykluskomplexes \(C_{\max}(K)\) zusammen.

\paragraph{Proposition 12.2 (Randzerstörung).}
Zum Kollapszeitpunkt \(t^{\ast}\) ist \(\partial\Omega\) keine wohldefinierte
trennende Oberfläche mehr zwischen zulässigen und unzulässigen Zuständen.
Formal gilt entweder
\begin{equation}
  \partial\Omega\big(K(t^{\ast})\big) = \emptyset
  \quad\text{oder}\quad
  \partial\Omega\big(K(t^{\ast})\big) = \Omega(M),
\end{equation}
je nachdem, ob die Einbettungsrestriktionen trivial streng oder trivial locker
werden. In beiden Fällen geht die strukturelle Rolle von \(\partial\Omega\)
als Viabilitätsgrenze verloren.

Die Operator-Perspektive lautet:
\begin{itemize}
  \item der Randoperator \(R\) ist undefiniert, sobald \(\Omega(K)=\emptyset\);
  \item der Continuumness-Operator \(U\) liefert \(k=0\) und bleibt danach für
        dieses Kontinuum konstant.
\end{itemize}

\subsection{Interner vs. externer Kollaps}

Kollaps kann durch innere Dynamik des Kontinuums oder durch äußere Veränderungen
seines Einbettungsraums ausgelöst werden. Diese Unterscheidung ist strukturell
nicht kausal: sie hängt davon ab, ob die Operatoren auf \(K\) oder auf \(M\) für
die Verletzung von Schwellwerten verantwortlich sind.

\paragraph{Definition 12.3 (Innerer Kollaps).}
\emph{Innerer Kollaps} liegt vor, wenn der Kollaps ausschließlich durch den
inneren Evolutionsoperator \(E=(F,G,H,Q,R,S,U)\) auf \(K\) ausgelöst wird,
während der Einbettungsraum \(M\) im relevanten Zeitintervall strukturell
statisch bleibt. Typische Signaturen sind:
\begin{itemize}
  \item destruierende Flüsse \(J_{\mathrm{kill}}\) dominieren stützende Flüsse;
  \item interne Schwellwerte \(\Theta_{\mathrm{stab}}\) oder
        \(\Theta_{\mathrm{crit}}\) werden wegen interner Spannungsakkumulation
        überschritten;
  \item Zykluszerfall, getrieben durch interne Imbalancen.
\end{itemize}

\paragraph{Definition 12.4 (Externer Kollaps).}
\emph{Externer Kollaps} liegt vor, wenn sich \(M\) derart ändert, dass keine
Konfiguration von \(K\) mit den neuen Restriktionen kompatibel bleibt.
Formal existiert ein Zeitintervall \([t_0,t^{\ast}]\), sodass:
\[
  \Omega(K(t_0))\neq\emptyset,
  \quad
  E\ \text{bleibt innerhalb des vorherigen zulässigen Bereichs},
\]
aber eine durch den eigenen Evolutionsoperator \(E_M\) induzierte Änderung von
\(M\) liefert
\[
  \Omega(K(t^{\ast})) =
    \big\{ s \in \Omega(M(t^{\ast})) \mid
           \Theta(K(t_0))(s)\le 0 \big\} = \emptyset.
\]

Folgerung~3.1 in Abschnitt~\ref{sec:results} kann wie folgt neu formuliert
werden: Externer Kollaps ist äquivalent zum Verlust gemeinsamer Lösungen der
Restriktionen von \(K\) und \(M\).

\paragraph{Gemischter Kollaps.}
In realen Systemen wirken oft beide Mechanismen zusammen. So kann eine
Protokolle intern Ressourcen verzehren, während eine externe Änderung der
osmotischen Bedingungen gleichzeitig die Viabilitätsbereiche verengt.
OC versucht nicht, diese Fälle hart zu trennen; die Unterscheidung intern/extern
ist vor allem ein Analysewerkzeug.

\paragraph{Kollapsdynamik}

Kollaps ist im Allgemeinen kein instantanes Ereignis auf der Ebene beobachtbarer
Größen. Der strukturelle Rahmen sagt eine charakteristische Annäherung an den
Kollaps vorher, gesteuert durch Schwellwerte, Spannung und Randgeometrie.

\paragraph{Kritische Verlangsamung.}
In der Nähe stabiler oder kritischer Schwellwerte \(\Theta_{\mathrm{stab}}\) und
\(\Theta_{\mathrm{crit}}\) relaxieren kleine Störungen zunehmend langsamer.
In der Operator-Sprache entspricht dies:
\begin{itemize}
  \item Realteile von Eigenwerten des linearen Flussoperators \(F\) nähern sich
        null;
  \item Zeitskalen der Zykluswiederherstellung, kodiert in \(Q\), divergieren;
  \item die Randbewegung unter \(R\) wird träge, bleibt aber auf Kontraktion von
        \(\Omega\) ausgerichtet.
\end{itemize}
Kritische Verlangsamung tritt bei physischen Phasenübergängen, ökologischen
Kollapsen, Finanzkrisen, Protokolleversagen und kognitiver Überlastung auf.

\paragraph{Divergenz der strukturellen Spannung.}
Die strukturelle Spannung \(T(K,t)\) nähert sich beim Kollaps dem Todesschwellwert
\(\Theta_{\mathrm{death}}(K)\):
\[
  \lim_{t\nearrow t^{\ast}} T\big(K,t\big)
    = \Theta_{\mathrm{death}}(K).
\]
Unterhalb dieses Schwellenwertes kann das System noch Zyklen erhalten, jedoch mit
immer enger werdenden Margen. Am Grenzwert zwingt jede weitere Störung
Trajektorien dazu, \(\partial\Omega\) zu überschreiten oder den
Einbettungsbereich zu verlassen.

\paragraph{Patch-Versagen und randgetriebener Kollaps.}
Das Patch-Modell der Grenzen aus
Abschnitt~\ref{sec:boundary-extended} liefert ein detailreicheres Bild.
Schreibt man den Rand als Sammlung von Patches \(\{P_i\}\) mit lokalen Zuständen
\(\sigma_i\) und Schwellwertvektoren \(\Theta_i\), kann Kollaps über
Folgendes verlaufen:
\begin{itemize}
  \item lokal verlierender Integritätsverlust (z.\,B. Membranruptur an einzelnen
        Patches, institutionelles Versagen in einzelnen Domänen);
  \item Weitergabe der Störung über Kopplung von Patches, die zu einem globalen
        Durchbruch des Randes und unkontrolliertem Austausch mit dem
        Einbettungsraum führt;
  \item schließlich Zerstörung aller Patches, die ein nichtleeres
        \(\Omega(K)\) stützen.
\end{itemize}
Aus Operator-Sicht handelt es sich dabei um eine gekoppelte Instabilität von
\(R\), \(H\) und \(Q\).

\paragraph{Zykluszerfall.}
Mit Annäherung an den Kollaps eliminiert der Zyklusoperator \(Q\) zunehmend
mehr Zyklen aus dem strukturstabilen Komplex. Die Folge
\[
  C_{\max}(t_0)
  \supset C_{\max}(t_1)
  \supset \dots
  \supset C_{\max}(t_n)
\]
verkleinert sich bis sie am Zeitpunkt \(t^{\ast}\) leer wird. Operativ:
\begin{itemize}
  \item Amplituden zentraler Zyklen nehmen ab;
  \item Zyklusperioden verlängern sich oder werden unregelmäßig;
  \item Zyklen tangieren immer häufiger \(\partial\Omega\) und werden durch
        Schwellwertüberschreitungen zerstört.
\end{itemize}

\subsection{Post-Kollaps-Rest}

Auch bei einem Kollaps kann der Kontinuumstrace im Einbettungsraum fortbestehen.
OC trennt strikt zwischen dem \emph{Kontinuum} und seinem
\emph{Residuum}.

\paragraph{Definition 12.5 (Residuum).}
Das \emph{Residuum} eines kollabierten Kontinuums \(K\) ist die Menge von
Strukturen im Einbettungsraum, die nach \(\Omega(K)=\emptyset\) verbleiben; Beispiele:
\begin{itemize}
  \item materielle Reste (z.\,B. Reaktionsprodukte, zerstörte Infrastruktur);
  \item topologische Reste (z.\,B. Defekte, Narben in Feldkonfigurationen);
  \item informationelle Reste (z.\,B. Dokumente, Datenspuren);
  \item normative Reste (z.\,B. teilweise erhaltene Rechtsrahmen).
\end{itemize}

Residuen sind nicht mehr durch das ursprüngliche Tupel
\((\Omega,A,P,J,\Theta,\partial\Omega,C,k)\) organisiert. Sie können jedoch das
Substrat für neue Kontinua bilden.

\paragraph{Wiederherstellungslimits.}
Da der Tod strukturell irreversibel ist (Satz~3 und Korollar~3.2 in
Abschnitt~\ref{sec:results}), kann kein Operator derselben Ebene den
Ausgangszustand aus dem Residuum rekonstruieren. Es werden zwei Fälle
unterschieden:
\begin{itemize}
  \item \emph{Rekonstruktion.}\
        Ein externer Akteur (z.\,B. Experimentator, Ingenieur, Institution)
        kann ein neues Kontinuum \(K'\) aufbauen, das in einzelnen Variablen
        \(K\) ähnelt; strukturell ist dies ein neues Geburtsereignis, keine
        Wiederherstellung.
  \item \emph{Spontane Reorganisation.}\
        Unter geeigneten Bedingungen können sich Residuuen zu einem neuen
        Kontinuum umorganisieren (siehe unten). Auch das ist eine neue
        Entität mit eigener Identität, selbst wenn viele Komponenten geerbt sind.
\end{itemize}

\subsection{Mechanismen der Wiedergeburt}

Wiedergeburt ist das Auftreten eines neuen Kontinuums nach dem Kollaps eines
vorherigen. OC behandelt Wiedergeburt als Spezialfall von Geburt, gesteuert durch
die gleichen Operatoren \(\Psi_{x\to x+1}\), aber mit Anfangsbedingungen aus
Residuuen.

\paragraph{Definition 12.6 (Wiedergeburt).}
Ein \emph{Wiedergeburtsevent} liegt vor, wenn nach dem Kollaps von \(K\) ein
neues Kontinuum \(K'\) im gleichen oder einem benachbarten
Einbettungsraum erscheint, sodass:
\begin{enumerate}
  \item es ein Zeitintervall gibt mit
        \(\Omega(K)=\emptyset\) und \(k(K,t)=0\), während
        \(\Omega(K')\neq\emptyset\);
  \item mindestens ein strukturelles Element von \(K'\) (Achsen, Potentiale,
        Randkomponenten, Residuen in \(\Omega(M)\)) vom Residuum von \(K\)
        übernommen wird;
  \item \(K'\) die Geburtsbedingungen auf einem Niveau \(K_y\) erfüllt, das
        dem ursprünglichen Niveau von \(K\) entsprechen kann oder ein anderes
        sein kann.
\end{enumerate}

Wiedergeburt ist deshalb strukturell selten und eingeschränkt. Die folgenden
Regime sind typisch.

\paragraph{Dimensionsbezogene Wiedergeburt.}
Nach dem Kollaps eines Kontinuums \(K_x\) können Residuuen die Geburt eines neuen
Kontinuums auf einem anderen Niveau tragen:
\begin{itemize}
  \item \emph{Aufwärts-Wiedergeburt} (\(K_x\to K_{x+1}\)):
        Kollaps einer überbeanspruchten Struktur auf Niveau \(K_x\) schafft
        Bedingungen für ein höherdimensionales Kontinuum (z.\,B. der
        Zusammenbruch eines einfachen institutionellen Rahmens, der eine
        komplexere Governance-Struktur ermöglicht);
  \item \emph{Seitwärts-Wiedergeburt} (\(K_x\to K_x'\)):
        Residuuen tragen ein neues Kontinuum auf demselben Niveau, aber mit
        anderen Achsen oder Schwellwerten (z.\,B. Reorganisation einer
        fehlgeschlagenen Protokolle zu einer neuen Protokolle mit
        veränderter Zusammensetzung).
\end{itemize}
Abwärts-Wiedergeburt (\(K_x\to K_{x-1}\)) gilt nicht als Wiedergeburt des
gleichen Kontinuums: Untere Ebenen können als unabhängige Entitäten bestehen
bleiben, doch das kollabierte höherstufige Kontinuum wird nicht
wiederbegründet.

\paragraph{Formale Bedingungen.}
Sei \(K\) zum Zeitpunkt \(t^{\ast}\) kollabiert und \(R_{\mathrm{res}}\) sein
Residuum in \(M\). Ein Wiedergeburtsevent ist möglich, wenn:
\begin{enumerate}
  \item das Residuum eine nichtleere Kandidatenzustandsmenge
        \(\Omega_{\mathrm{res}}\subseteq \Omega(M)\) definiert;
  \item Achsen \(A_{\mathrm{res}}\subseteq A(M)\) und Schwellwerte
        \(\Theta_{\mathrm{res}}\) existieren mit
        \[
          \Omega(K') =
            \{ s\in \Omega_{\mathrm{res}}
               \mid \Theta_{\mathrm{res}}(s)\le 0
            \}
          \neq \emptyset;
        \]
  \item die strukturelle Spannung auf Basis von
        \((\Omega_{\mathrm{res}},A_{\mathrm{res}},\Theta_{\mathrm{res}})\)
        den zugehörigen Dimensionsschwellwert überschreitet (falls ein
        Niveauanstieg stattfindet) und ein passender Geburtsoperator
        \(\Psi\) definiert ist.
\end{enumerate}

\paragraph{Beispiele über \texorpdfstring{\(K_3\)}{K_3}–\texorpdfstring{\(K_5\)}{K_5}.}
Im chemischen und protocellulären Regime:
\begin{itemize}
  \item der Kollaps einer Protokolle (Membranruptur, Gradientenverlust) lässt
        Reste aus Lipiden, Nukleotiden und Katalysatoren in der Umgebung;
  \item diese Reste können unter günstigen Bedingungen in neue Vesikel und
        Netzwerke re-kombinieren;
  \item die neuen Protokollen bilden neue Kontinua \(K_4'\) mit eigenen
        Schwellwerten und Zyklen, selbst wenn sie Komponenten der kollabierten
        Vorläufer übernehmen.
\end{itemize}

In frühen bioelektrischen Systemen:
\begin{itemize}
  \item Kollaps der Erregbarkeit in einem primitiven Netzwerk (z.\,B. wegen
        Kanalausfall) kann Membranfragmente und Kanäle hinterlassen, die neue
        erregbare Patches bilden;
  \item regenerierte erregbare Domänen definieren neue \(K_5'\)-Kontinua.
\end{itemize}

\paragraph{Wiedergeburt auf höheren Ebenen.}
Auf kognitiver, sozialer und zivilisatorischer Ebene:
\begin{itemize}
  \item kognitiver Kollaps (Verlust kohärenter interner Modelle) kann von der
        Wiederaufbauung neuer Modelle mit intakten Repräsentationen als
        Ausgangspunkt folgen; strukturell ist dies die Geburt eines neuen
        kognitiven Kontinuums \(K_6'\);
  \item institutioneller Kollaps (z.\,B. Scheitern eines Regierungsregimes)
        hinterlässt rechtliche, organisatorische und infrastrukturelle Reste,
        die für die Bildung neuer Institutionen \(K_7'\) verwendet werden können;
  \item zivilisatorischer Kollaps hinterlässt materielle Infrastruktur,
        kulturelle Artefakte und Bevölkerungsverteilungen, die emergente
        Zivilisationen \(K_8'\) tragen können.
\end{itemize}

In allen Fällen fordert OC, dass das neue Kontinuum eine eigene Identität hat; die
strukturelle Irreversibilität des Todes bleibt unberührt.

\subsection{Kollaps in höherstufigen Kontinua}

Höhere Ebenen \(K_6\)–\(K_{10}\) zeigen Kollapsphänomene, die weniger
offensichtlich physisch sind, aber strukturell analog wirken.

\paragraph{Kognitiver Kollaps (\texorpdfstring{\(K_6\)}{K_6}).}
Für kognitive Kontinua sind die Kernelemente:
\begin{itemize}
  \item Achsen für Konzepte, Merkmale und Modellparameter;
  \item Potentiale, die Vorhersagefehler, Wert und Vertrauen messen;
  \item Zyklen, die Aufmerksamkeits-Schleifen, Vorhersage--Korrektur-Zyklen
        und Lernzyklen repräsentieren.
\end{itemize}
Kognitiver Kollaps liegt vor, wenn:
\begin{itemize}
  \item Vorhersagefehler und innere Spannung die Ausdrücklichkeitsschwellen
        \(\Theta_{\mathrm{expr}}\) überschreiten;
  \item kein kohärentes internes Modell innerhalb gegebener Achsen und
        Schwellwerte gehalten werden kann;
  \item alle strukturstabilen kognitiven Zyklen verschwinden und
        \(C_{\max}(K_6)=\emptyset\) sowie \(k(K_6,t)\to 0\) gilt.
\end{itemize}
Beispiele sind vollständiger Zusammenbruch eines konzeptuellen Schemas, starke
Überlastung oder pathologische Zustände, in denen das System kein stabiles Modell
seiner Umwelt aufrechterhalten kann.

\paragraph{Institutioneller Kollaps (\texorpdfstring{\(K_7\)}{K_7}).}
Für soziale Kontinua sind die relevanten Komponenten:
\begin{itemize}
  \item soziale Graphen, Rollenstrukturen und institutionelle Anordnungen in
        \(\Omega(K_7)\);
  \item Vertrauen, Legitimität und Ressourcenschwellwerte in \(\Theta(K_7)\);
  \item Zyklen, die institutionelle Routinen und Governance-Schleifen
        repräsentieren.
\end{itemize}
Institutioneller Kollaps liegt vor, wenn:
\begin{itemize}
  \item Vertrauen und Legitimität unter Existenzschwellwerte sinken;
  \item zentrale Zyklen (z.\,B. Steuererhebung, Rechtsdurchsetzung,
        Entscheidungs-Loop) brechen;
  \item soziale Graphen stärker fragmentieren, als mit institutioneller
        Funktionsweise vereinbar ist.
\end{itemize}
Das Residuum kann juristische Texte, physische Gebäude und fragmentierte
Netzwerke umfassen, aus denen später neue Institutionen \(K_7'\) aufgebaut
werden können.

\paragraph{Zivilisatorischer Kollaps (\texorpdfstring{\(K_8\)}{K_8}).}
Auf der Ebene der Zivilisation:
\begin{itemize}
  \item Achsen repräsentieren Infrastrukturen, Sektoren, Regionen und
        großskalige Organisationsformen;
  \item Potentiale umfassen Energieverfügbarkeit, Ressourcenbestände und
        Risikogradienten;
  \item Zyklen beinhalten Wirtschaftszyklen, Erneuerung von Infrastruktur,
        Wissensweitergabe.
\end{itemize}
Zivilisatorischer Kollaps entspricht einem Zusammenbruch dieser Zyklen über
mehrere Sektoren hinweg, dem Überschreiten mehrerer Schwellwerte der
erweiterten Schwellwertlandschaft und der globalen Kontraktion von
\(\Omega(K_8)\) auf die leere Menge. Das Residuum besteht aus
Infrastruktur, Bevölkerungsverteilungen und kulturellen Artefakten, die künftige
Zivilisationen hervorrufen können.

\paragraph{Theoretischer und meta-theoretischer Kollaps (\texorpdfstring{\(K_9\)}{K_9}–\texorpdfstring{\(K_{10}\)}{K_{10}}).}
Auch Theorien und Meta-Theorien können strukturell kollabieren:
\begin{itemize}
  \item innere Inkonsistenz oder empirische Fehlgriffe können
        \(\Theta_{\mathrm{exist}}\) oder \(\Theta_{\mathrm{stab}}\) auf
        Niveau \(K_9\) verletzen;
  \item Meta-theoretische Rahmen auf \(K_{10}\) können inkohärent werden oder
        keine wachsende Modellmenge integrieren können, wodurch
        Ausdrucksschwellen verletzt werden;
  \item in beiden Fällen verschwinden Zykluskomplexe, die Forschungsprogramme,
        Paradigmenzyklen und Meta-Update-Schleifen repräsentieren, bis
        \(C_{\max}=\emptyset\) und das Framework de facto stirbt.
\end{itemize}
Das anschließende Auftreten neuer Theorien oder Meta-Rahmenwerke, die auf
überlebenden Ergebnissen aufbauen, wird strukturell als Wiedergeburt
behandelt.

\subsection{Zusammenfassung}

Kollaps und Wiedergeburt in OC sind keine Metaphern, sondern strukturell
definierte Phänomene. Kollaps entspricht dem Leerenwerden des zulässigen
Zustandsraums, dem Verschwinden stabiler Zyklen, dem Verlust einer sinnvollen
Grenze und dem Verschwinden der Continuumness. Er kann durch interne Dynamik,
durch externe Änderungen im Einbettungsraum oder durch deren Kopplung ausgelöst
werden und zeigt charakteristische Signaturen wie kritische Verlangsamung,
Divergenz struktureller Spannung, Patch-Versagen und Zykluszerfall.

Wiedergeburt ist stark eingeschränkt: Ein neues Kontinuum kann aus Residuuen nur
entstehen, wenn Einbettungsräume und Schwellwertlandschaften wieder eine
nichtleere zulässige Region tragen und die Geburtsbedingungen erfüllt sind. Das
neue Kontinuum besitzt stets eine eigene strukturelle Identität, auch wenn es
Komponenten des kollabierten Vorgängers übernimmt.

Diese Prinzipien gelten einheitlich vom frühen Protokolle und frühen
bioelektrischen Systemen bis zu kognitiven und sozialen Strukturen sowie
Theorien und meta-theoretischen Rahmenwerken. Kollaps und Wiedergeburt sind
damit keine Sonderfälle, sondern generische Regime derselben strukturellen
Mechanik, die Geburt, Evolution und Tod von Kontinua über die Hierarchie
\(K_0\)–\(K_{10}\) steuert.




